\documentclass[12pt,a4paper]{article}
\input{../../common/header}

\setlist{nosep}

\begin{document}

Написати функцію $f$, яка отримує

\begin{enumerate}
    \item дійсні числа $x$ та $y$ і повертає логічну ознаку того, що точка $(x, y)$ належить області визначення виразу $$\frac{1}{x} + \frac{1}{y}$$
    \item дійсні числа $x$ та $y$ і повертає логічну ознаку того, що точка $(x, y)$ належить області визначення виразу $$\frac{1}{x \cdot y}$$
    \item дійсні числа $x$ та $y$ і повертає логічну ознаку того, що точка $(x, y)$ належить області визначення виразу $$\frac{1}{x + y}$$
    \item дійсні числа $x$ та $y$ і повертає логічну ознаку того, що точка $(x, y)$ належить області визначення виразу $$\frac{2025}{x + y}$$
    \item дійсні числа $x$ та $y$ і повертає логічну ознаку того, що точка $(x, y)$ належить області визначення виразу $$\frac{1}{\sqrt{x}} + \frac{1}{\sqrt{y}}$$
    \item дійсні числа $x$ та $y$ і повертає логічну ознаку того, що точка $(x, y)$ належить області визначення виразу $$\frac{1}{\sqrt{x \cdot y}}$$
    \item дійсні числа $x$ та $y$ і повертає логічну ознаку того, що точка $(x, y)$ належить області визначення виразу $$\frac{2024}{\sqrt{x \cdot y}}$$
    \item дійсні числа $x$ та $y$ і повертає логічну ознаку того, що точка $(x, y)$ належить області визначення виразу $$\frac{1}{\sqrt{x + y}}$$

    \item дійсні числа $x$ та $y$ і повертає логічну ознаку того, що точка $(x, y)$ належить області визначення виразу $$\frac{1}{x+1} + \frac{1}{y-1}$$
    \item дійсні числа $x$ та $y$ і повертає логічну ознаку того, що точка $(x, y)$ належить області визначення виразу $$\frac{1}{(x-1) \cdot (y+1)}$$
    \item дійсні числа $x$ та $y$ і повертає логічну ознаку того, що точка $(x, y)$ належить області визначення виразу $$\frac{1}{x + y + 1}$$
    \item дійсні числа $x$ та $y$ і повертає логічну ознаку того, що точка $(x, y)$ належить області визначення виразу $$\frac{1}{\sqrt{x-1}} + \frac{1}{\sqrt{y+1}}$$
    \item дійсні числа $x$ та $y$ і повертає логічну ознаку того, що точка $(x, y)$ належить області визначення виразу $$\frac{1}{\sqrt{x + y - 1}}$$
    
    \item дійсні числа $x$ та $y$ і повертає логічну ознаку того, що точка $(x, y)$ належить області визначення виразу $$\frac{1}{x} + \sqrt{y}$$
    \item дійсні числа $x$ та $y$ і повертає логічну ознаку того, що точка $(x, y)$ належить області визначення виразу $$\frac{1}{x} \cdot \sqrt{y}$$
    \item дійсні числа $x$ та $y$ і повертає логічну ознаку того, що точка $(x, y)$ належить області визначення виразу $$\frac{1}{x} + y$$
    \item дійсні числа $x$ та $y$ і повертає логічну ознаку того, що точка $(x, y)$ належить області визначення виразу $$x + \frac{1}{\sqrt{y}}$$
    \item дійсні числа $x$ та $y$ і повертає логічну ознаку того, що точка $(x, y)$ належить області визначення виразу $$\frac{\sqrt{x}}{\sqrt{y}}$$
    \item дійсні числа $x$ та $y$ і повертає логічну ознаку того, що точка $(x, y)$ належить області визначення виразу $$\frac{y}{\sqrt{x}}$$

    \item дійсні числа $a$, $b$ та $c$ і повертає логічну ознаку того, що точка $(a, b, c)$ належить області визначення виразу $$\frac{1}{a} + \frac{1}{b} + c$$
    \item дійсні числа $a$, $b$ та $c$ і повертає логічну ознаку того, що точка $(a, b, c)$ належить області визначення виразу $$\frac{c}{a \cdot b}$$
    \item дійсні числа $a$, $b$ та $c$ і повертає логічну ознаку того, що точка $(a, b, c)$ належить області визначення виразу $$\frac{c}{a + b}$$
    \item дійсні числа $a$, $b$ та $c$ і повертає логічну ознаку того, що точка $(a, b, c)$ належить області визначення виразу $$\frac{1}{\sqrt{a}} + \frac{1}{\sqrt{b}} + c$$
    \item дійсні числа $a$, $b$ та $c$ і повертає логічну ознаку того, що точка $(a, b, c)$ належить області визначення виразу $$\frac{1}{\sqrt{a}} + \frac{b}{\sqrt{c}}$$
    \item дійсні числа $a$, $b$ та $c$ і повертає логічну ознаку того, що точка $(a, b, c)$ належить області визначення виразу $$\frac{c}{\sqrt{a \cdot b}}$$
    \item дійсні числа $a$, $b$ та $c$ і повертає логічну ознаку того, що точка $(a, b, c)$ належить області визначення виразу $$\frac{1}{\sqrt{a + b}} + c$$
    \item дійсні числа $a$, $b$ та $c$ і повертає логічну ознаку того, що точка $(a, b, c)$ належить області визначення виразу $$\frac{2024}{\sqrt{a + b}} + c$$

    \item цілі числа $n$ та $m$ і повертає логічну ознаку того, що точка $(n, m)$ належить області визначення виразу $$\frac{1}{n} + \frac{1}{m}$$
    \item цілі числа $n$ та $m$ і повертає логічну ознаку того, що точка $(n, m)$ належить області визначення виразу $$\frac{1}{n \cdot m}$$
    \item цілі числа $n$ та $m$ і повертає логічну ознаку того, що точка $(n, m)$ належить області визначення виразу $$\frac{1}{n + m}$$
    \item цілі числа $n$ та $m$ і повертає логічну ознаку того, що точка $(n, m)$ належить області визначення виразу $$\frac{2024}{n + m}$$
    \item цілі числа $n$ та $m$ і повертає логічну ознаку того, що точка $(n, m)$ належить області визначення виразу $$\frac{1}{\sqrt{n}} + \frac{1}{\sqrt{m}}$$
    \item цілі числа $n$ та $m$ і повертає логічну ознаку того, що точка $(n, m)$ належить області визначення виразу $$\frac{1}{\sqrt{n \cdot m}}$$
    \item цілі числа $n$ та $m$ і повертає логічну ознаку того, що точка $(n, m)$ належить області визначення виразу $$\frac{2025}{\sqrt{n \cdot m}}$$
    \item цілі числа $n$ та $m$ і повертає логічну ознаку того, що точка $(n, m)$ належить області визначення виразу $$\frac{1}{\sqrt{n + m}}$$
    
    \item цілі числа $n$ та $m$ і повертає логічну ознаку того, що точка $(n, m)$ належить області визначення виразу $$\frac{1}{n+1} + \frac{1}{m-1}$$
    \item цілі числа $n$ та $m$ і повертає логічну ознаку того, що точка $(n, m)$ належить області визначення виразу $$\frac{1}{(n-1) \cdot (m+1)}$$
    \item цілі числа $n$ та $m$ і повертає логічну ознаку того, що точка $(n, m)$ належить області визначення виразу $$\frac{1}{n + m + 1}$$
    \item цілі числа $n$ та $m$ і повертає логічну ознаку того, що точка $(n, m)$ належить області визначення виразу $$\frac{1}{\sqrt{n-1}} + \frac{1}{\sqrt{m+1}}$$
    \item цілі числа $n$ та $m$ і повертає логічну ознаку того, що точка $(n, m)$ належить області визначення виразу $$\frac{1}{\sqrt{n + m - 1}}$$
    
    \item цілі числа $n$ та $m$ і повертає логічну ознаку того, що точка $(n, m)$ належить області визначення виразу $$\frac{1}{n} + \sqrt{m}$$
    \item цілі числа $n$ та $m$ і повертає логічну ознаку того, що точка $(n, m)$ належить області визначення виразу $$\frac{1}{n} \cdot \sqrt{m}$$
    \item цілі числа $n$ та $m$ і повертає логічну ознаку того, що точка $(n, m)$ належить області визначення виразу $$\frac{1}{n} + m$$
    \item цілі числа $n$ та $m$ і повертає логічну ознаку того, що точка $(n, m)$ належить області визначення виразу $$n + \frac{1}{\sqrt{m}}$$
    \item цілі числа $n$ та $m$ і повертає логічну ознаку того, що точка $(n, m)$ належить області визначення виразу $$\frac{\sqrt{n}}{\sqrt{m}}$$
    \item цілі числа $n$ та $m$ і повертає логічну ознаку того, що точка $(n, m)$ належить області визначення виразу $$\frac{m}{\sqrt{n}}$$
    
    \item цілі числа $n$, $m$ та $k$ і повертає логічну ознаку того, що точка $(n, m, k)$ належить області визначення виразу $$\frac{1}{n} + \frac{1}{m} + k$$
    \item цілі числа $n$, $m$ та $k$ і повертає логічну ознаку того, що точка $(n, m, k)$ належить області визначення виразу $$\frac{k}{n \cdot m}$$
    \item цілі числа $n$, $m$ та $k$ і повертає логічну ознаку того, що точка $(n, m, k)$ належить області визначення виразу $$\frac{k}{n + m}$$
    \item цілі числа $n$, $m$ та $k$ і повертає логічну ознаку того, що точка $(n, m, k)$ належить області визначення виразу $$\frac{1}{\sqrt{n}} + \frac{1}{\sqrt{m}} + k$$
    \item цілі числа $n$, $m$ та $k$ і повертає логічну ознаку того, що точка $(n, m, k)$ належить області визначення виразу $$\frac{1}{\sqrt{n}} + \frac{m}{\sqrt{k}}$$
    \item цілі числа $n$, $m$ та $k$ і повертає логічну ознаку того, що точка $(n, m, k)$ належить області визначення виразу $$\frac{k}{\sqrt{n \cdot m}}$$
    \item цілі числа $n$, $m$ та $k$ і повертає логічну ознаку того, що точка $(n, m, k)$ належить області визначення виразу $$\frac{1}{\sqrt{n + m}} + k$$
    \item цілі числа $n$, $m$ та $k$ і повертає логічну ознаку того, що точка $(n, m, k)$ належить області визначення виразу $$\frac{2025}{\sqrt{n + m}} + k$$
\end{enumerate}

Стиль запису умови також оцінюється.

\end{document}